\documentclass[lang=cn,12pt]{elegantbook}

\title{数分高代习题课讲义}
\subtitle{大一年级下学期 }

\author{忘记了一天}
\institute{XTU\_data\_2}
\date{April 2, 2026}
% \version{4.3}
% \bioinfo{自定义}{信息}

\extrainfo{世界上只有一种真正的英雄主义，那就是在认清生活真相之后依然热爱生活。}

\setcounter{tocdepth}{3}

\logo{XTU_logo_1.jpg}
\cover{100.jpg}

% 本文档命令
\usepackage{array}
\newcommand{\ccr}[1]{\makecell{{\color{#1}\rule{1cm}{1cm}}}}

% 修改标题页的橙色带
\definecolor{customcolor}{RGB}{154, 206, 246}
\colorlet{coverlinecolor}{customcolor}

\begin{document}

\maketitle

\frontmatter

\tableofcontents

\mainmatter

% \chapter*{综述}
% \markboth{Introduction}{Introduction}


\chapter{第一次习题课}

\section{线性映射相关知识}
\begin{definition}{线性映射}
    设 \(V\) 和 \(U\) 是数域 \(K\) 上的线性空间，\(\varphi\) 是 \(V\) 到 \(U\) 的映射，如果对任意 \(\alpha, \beta \in V\) 和任意数 \(k\)，都有
    \[
    \varphi(\alpha + \beta) = \varphi(\alpha) + \varphi(\beta), \quad \varphi(k\alpha) = k\varphi(\alpha),
    \]
    则称 \(\varphi\) 是 \(V\) 到 \(U\) 的线性映射。
\end{definition}

\begin{definition}{线性变换}
    设 \(V\) 是数域 \(K\) 上的线性空间，\(\varphi\) 是 \(V\) 到 \(V\) 的映射，如果 \(\varphi\) 是 \(V\) 到 \(V\) 的线性映射，则称 \(\varphi\) 是 \(V\) 上的线性变换。
\end{definition}

\begin{definition}{像空间}
    设 \(\varphi\) 是线性空间 \(V\) 到 \(U\) 的线性映射，\(\varphi(V) = \{\varphi(\alpha): \alpha \in V\}\)，称 \(\varphi(V)\) 是线性映射 \(\varphi\) 的像空间，记作 \(Im\varphi\)。
\end{definition}

\begin{definition}{核空间}
    设 \(\varphi\) 是线性空间 \(V\) 到 \(U\) 的线性映射，\(Ker\varphi = \{\alpha \in V: \varphi(\alpha) = 0\}\)，称 \(Ker\varphi\) 是线性映射 \(\varphi\) 的核空间，记作 \(Ker\varphi\)。
\end{definition}

\begin{theorem}
    设 \(\varphi\) 是数域 \(\mathbb{F}\) 上向量空间 \(V\) 到 \(U\) 的线性映射，\(A\) 是 \(\varphi\) 在任意给定基下的表示矩阵，则 \(\varphi\) 的秩等于 \(A\) 的秩，\(\varphi\) 的零度等于 \(V\) 的维数减去 \(A\) 的秩。即有
\[
\dim \operatorname{Im} \varphi + \dim \operatorname{Ker} \varphi = \dim V.
\]
\end{theorem}

\begin{definition}{不变子空间}
    设 \(\varphi\) 是线性空间 \(V\) 上的线性变换，如果 \(V\) 的子空间 \(W\) 满足 \(\varphi(W) \subseteq W\)，则称 \(W\) 是线性变换 \(\varphi\) 的不变子空间。
\end{definition}

\begin{theorem}
    设 \(\varphi\) 是数域 \(\mathbb{F}\) 上向量空间 \(V\) 上的线性变换，\(W\) 是 \(\varphi\) 的不变子空间。若取 \(W\) 的一组基 \(\{e_1, \cdots, e_r\}\)，再扩张为 \(V\) 的一组基 \(\{e_1, \cdots, e_r, e_{r+1}, \cdots, e_n\}\)，则 \(\varphi\) 在这组基下的表示矩阵具有下列分块上三角矩阵的形状：
\[
\begin{pmatrix}
A_{11} & A_{12} \\
O & A_{22}
\end{pmatrix},
\]
其中 \(A_{11}\) 是一个 \(r\) 阶矩阵。
\end{theorem}


\section{习题讲解}

\begin{example}\label{ex4.8}
    设 \(\varphi\) 是 \(n\) 维线性空间 \(V\) 上的线性变换，\(\alpha \in V\)。若 \(\varphi^{m-1}(\alpha) \neq 0\)，而  
    \(\varphi^m(\alpha) = 0\)，求证：\(\alpha, \varphi(\alpha), \varphi^2(\alpha), \cdots, \varphi^{m-1}(\alpha)\) 线性无关。
\end{example}
    
\begin{proof}
    设有 \(m\) 个数 \(a_0, a_1, \cdots, a_{m-1}\)，使得  
    \[
    a_0\alpha + a_1\varphi(\alpha) + \cdots + a_{m-1}\varphi^{m-1}(\alpha) = 0.
    \]
    
    上式两边同时作用 \(\varphi^{m-1}\)，则有  
    \[
    a_0\varphi^{m-1}(\alpha) = 0,
    \]
    由于 \(\varphi^{m-1}(\alpha) \neq 0\)，故 \(a_0 = 0\)。上式两边同时作用 \(\varphi^{m-2}\)，则有  
    \[
    a_1\varphi^{m-1}(\alpha) = 0,
    \]
    由于 \(\varphi^{m-1}(\alpha) \neq 0\)，故 \(a_1 = 0\)。不断这样做下去，最后可得  
    \[
    a_0 = a_1 = \cdots = a_{m-1} = 0,
    \]
    于是 \(\alpha, \varphi(\alpha), \varphi^2(\alpha), \cdots, \varphi^{m-1}(\alpha)\) 线性无关。
\end{proof}

\begin{example}
    设 \( V \) 是数域 \( K \) 上的 \( n \) 维线性空间，\( \varphi \) 是 \( V \) 上的幂零线性变换，满足  
    \( r(\varphi) = n - 1 \)。求证：存在 \( V \) 的一组基，使得 \( \varphi \) 在这组基下的表示矩阵为  
    \[
    A = 
    \begin{pmatrix}
    0 & 0 & \cdots & 0 \\
    1 & 0 & \cdots & 0 \\
    0 & 1 & \cdots & 0 \\
    \vdots & \vdots & \ddots & \vdots \\
    0 & 0 & \cdots & 1
    \end{pmatrix}.
    \]
\end{example}
    
\begin{proof}
    由假设存在正整数 \( m \)，使得 \( \varphi^m = 0, \varphi^{m-1} \neq 0 \)，从而存在 \( \alpha \in V \)，使得  
    \[
    \varphi^m(\alpha) = 0, \quad \varphi^{m-1}(\alpha) \neq 0。
    \]  
    由例\ref{ex4.8}可知，\( \alpha, \varphi(\alpha), \cdots, \varphi^{m-1}(\alpha) \) 线性无关，于是  
    \[
    m \leq \dim V = n。
    \]  
    另一方面，由 Sylvester 不等式$r(AB) \geq r(A) + r(B) -n $以及  
    \[
    r(\varphi) = n - 1
    \]  
    可知，
    \[
    r(\varphi^2) \geq 2r(\varphi) - n = n - 2。
    \]  
    不断这样讨论下去，最终可得  
    \[
    0 = r(\varphi^m) \geq n - m,
    \]  
    即有  
    \[
    m \geq n,
    \]  
    从而  
    \[
    m = n。
    \]  
    于是 \( \alpha, \varphi(\alpha), \cdots, \varphi^{n-1}(\alpha) \) 是 \( V \) 的一组基，\( \varphi \) 在这组基下的表示矩阵为 \( A \)。
\end{proof}

\begin{example}
    设 \(\varphi\) 是 \(n\) 维线性空间 \(V\) 上的线性变换，证明：\(\varphi\) 是可逆变换的充要条件是 \(\varphi\) 将 \(V\) 的基变为基。
\end{example}
    
\begin{proof}
若 \(\varphi\) 是可逆变换，则显然 \(\varphi\) 将 \(V\) 的基变为基。反之，若 \(e_1, e_2, \cdots, e_n\) 和 \(f_1, f_2, \cdots, f_n\) 是 \(V\) 的两组基，使得 \(\varphi(e_i) = f_i \ (1 \leq i \leq n)\)，则对任意 \(\alpha \in V\)，
\[
\alpha = \lambda_1 f_1 + \lambda_2 f_2 + \cdots + \lambda_n f_n,
\]
有 \(\varphi(\lambda_1 e_1 + \lambda_2 e_2 + \cdots + \lambda_n e_n) = \alpha\)，即 \(\varphi\) 是满映射，从而是自同构。我们也可以证明 \(\varphi\) 是单映射，从而是自同构（留给读者完成）。另外还可以这样讨论，设从基 \(e_1, e_2, \cdots, e_n\) 到基 \(f_1, f_2, \cdots, f_n\) 的过渡矩阵为 \(P\)，则 \(\varphi\) 在基 \(e_1, e_2, \cdots, e_n\) 下的表示矩阵就是 \(P\)，这是一个可逆矩阵，从而 \(\varphi\) 是可逆变换。
\end{proof}

\begin{example}\label{ex4.21}
    设 \(\varphi\) 是线性空间 \(V\) 到 \(U\) 的线性映射，\(\{e_1, e_2, \cdots, e_n\}\) 和 \(\{f_1, f_2, \cdots, f_n\}\) 是 \(V\) 的两组基，\(\{e_1, e_2, \cdots, e_n\}\) 到 \(\{f_1, f_2, \cdots, f_n\}\) 的过渡矩阵为 \(P\)。\(\{g_1, g_2, \cdots, g_m\}\) 和 \(\{h_1, h_2, \cdots, h_m\}\) 是 \(U\) 的两组基，\(\{g_1, g_2, \cdots, g_m\}\) 到 \(\{h_1, h_2, \cdots, h_m\}\) 的过渡矩阵为 \(Q\)。又设 \(\varphi\) 在基 \(\{e_1, e_2, \cdots, e_n\}\) 和基 \(\{g_1, g_2, \cdots, g_m\}\) 下的表示矩阵为 \(A\)，在基 \(\{f_1, f_2, \cdots, f_n\}\) 和基 \(\{h_1, h_2, \cdots, h_m\}\) 下的表示矩阵为 \(B\)。求证：
    \[
    B = Q^{-1}AP.
    \]
\end{example}
    
\begin{proof}
    任取 \(v \in V\)，设它在基 \(\{e_1, e_2, \cdots, e_n\}\) 下的坐标向量为 \((x_1, x_2, \cdots, x_n)'\)，则它在基 \(\{f_1, f_2, \cdots, f_n\}\) 下的坐标向量为 \(P^{-1}(x_1, x_2, \cdots, x_n)'\)。\(\varphi(v)\) 在基 \(\{g_1, g_2, \cdots, g_m\}\) 下的坐标向量为 \(A(x_1, x_2, \cdots, x_n)'\)，在基 \(\{h_1, h_2, \cdots, h_m\}\) 下的坐标向量为 \(BP^{-1}(x_1, x_2, \cdots, x_n)'\)。由于从 \(\{g_1, g_2, \cdots, g_m\}\) 到 \(\{h_1, h_2, \cdots, h_m\}\) 的过渡矩阵
    为 \(Q\)，故 \(A(x_1, x_2, \cdots, x_n)' = QBP^{-1}(x_1, x_2, \cdots, x_n)'\)。因为 \((x_1, x_2, \cdots, x_n)'\) 是任意的，故 \(A = QBP^{-1}\)，即 \(B = Q^{-1}AP\)。
\end{proof}


\begin{example}\label{ex4.22}
    设 \(\varphi\) 是有限维线性空间 \(V\) 到 \(U\) 的线性映射，求证：必存在 \(V\) 和 \(U\) 的两组基，使线性映射 \(\varphi\) 在两组基下的表示矩阵为  
    \[
    \begin{pmatrix} 
    I_r & O \\ 
    O & O 
    \end{pmatrix}.
    \]
\end{example}
    
\begin{proof}
    设 \(\{e_1, e_2, \cdots, e_n\}\) 是 \(V\) 的一组基，\(\{g_1, g_2, \cdots, g_m\}\) 是 \(U\) 的一组基，\(\varphi\) 在这两组基下的表示矩阵为 \(A\)。由相抵标准型理论可知，存在 \(m\) 阶非异阵 \(Q, n\) 阶非异阵 \(P\)，使得  
    \[
    Q^{-1}AP = 
    \begin{pmatrix} 
    I_r & O \\ 
    O & O 
    \end{pmatrix}.
    \]  
    设 \(\{f_1, f_2, \cdots, f_n\}\) 是 \(V\) 的一组新基，使得从 \(\{e_1, e_2, \cdots, e_n\}\) 到 \(\{f_1, f_2, \cdots, f_n\}\) 的过渡矩阵为 \(P\)；设 \(\{h_1, h_2, \cdots, h_m\}\) 是 \(U\) 的一组新基，使得从 \(\{g_1, g_2, \cdots, g_m\}\) 到 \(\{h_1, h_2, \cdots, h_m\}\) 的过渡矩阵为 \(Q\)，则由例 \ref{ex4.21} 可知，\(\varphi\) 在两组新基下的表示矩阵为  
    \[
    Q^{-1}AP = 
    \begin{pmatrix} 
    I_r & O \\ 
    O & O 
    \end{pmatrix}.
    \]
\end{proof}

% \begin{example}
%     设 \(\varphi\) 是有限维线性空间 \(V\) 到 \(U\) 的线性映射，求证：必存在 \(U\) 到 \(V\) 的线性映射 \(\psi\)，使得 \(\varphi\psi\varphi = \varphi\)。
% \end{example}
    
% \begin{proof}[证法 1（几何方法）]
%     设 \(V\) 和 \(U\) 的维数分别是 \(n\) 和 \(m\)。由例 \ref{ex4.22} 可知，存在 \(V\) 和 \(U\) 的基 \(\{e_1, e_2, \cdots, e_n\}, \{f_1, f_2, \cdots, f_m\}\)，使得 \(\varphi\) 在这两组基下的表示矩阵为
%     \[
%     \begin{pmatrix}
%     I_r & O \\
%     O & O
%     \end{pmatrix}.
%     \]
%     这就是 \(\varphi(e_i) = f_i, 1 \leq i \leq r; \varphi(e_j) = 0, r+1 \leq j \leq m\)。定义 \(\psi\) 是 \(U\) 到 \(V\) 的线性映射，它在基上的作用为
%     \[
%     \psi(f_i) = e_i, 1 \leq i \leq r; \quad \psi(f_j) = 0, r+1 \leq j \leq m,
%     \]
%     则在 \(V\) 的基上，有
%     \[
%     \varphi\psi\varphi(e_i) = \varphi\psi(f_i) = \varphi(e_i), \quad 1 \leq i \leq r;
%     \]
%     \[
%     \varphi\psi\varphi(e_j) = \varphi\psi(0) = 0 = \varphi(e_j), \quad r+1 \leq j \leq m.
%     \]
%     于是 \(\varphi\psi\varphi = \varphi\)。
% \end{proof}

\begin{example}
    设 \(\varphi\) 是线性空间 \(V\) 上的线性变换，若它在 \(V\) 的任一组基下的表示矩阵都相同，求证：\(\varphi\) 是纯量变换，即存在常数 \(k\)，使得 \(\varphi(\alpha) = k\alpha\) 对一切 \(\alpha \in V\) 都成立。
\end{example}
    
\begin{proof}
    取定 \(V\) 的一组基，设 \(\varphi\) 在这组基下的表示矩阵是 \(A\)。由已知条件可知，对任意一个同阶可逆矩阵 \(P\)，\(A = P^{-1}AP\)，即 \(PA = AP\)。因此矩阵 \(A\) 和任意一个可逆矩阵乘法可交换，于是 \(A = kI_n\)，由此即知 \(\varphi\) 是纯量变换。
\end{proof}

% \begin{example}\label{ex3.93}
%     设 \( A \) 为 \( m \times n \) 矩阵，证明：存在 \( n \times m \) 矩阵 \( B \)，使得 \( ABA = A \)。
% \end{example}
    
% \begin{proof}[证法 1]
%     设 \( PAQ = \begin{pmatrix} I_r & O \\ O & O \end{pmatrix} \)，其中 \( P \) 是 \( m \) 阶非异阵，\( Q \) 是 \( n \) 阶非异阵。注意到问题的条件和结论在相抵变换：\( A \mapsto PAQ \)，\( B \mapsto Q^{-1}BP^{-1} \) 下保持不变，故不妨从一开始就假设 \( A = \begin{pmatrix} I_r & O \\ O & O \end{pmatrix} \) 是相抵标准型。设 \( B = \begin{pmatrix} B_1 & B_2 \\ B_3 & B_4 \end{pmatrix} \) 为对应的分块，由 \( ABA = A \) 可得 \( B_1 = I_r \)，其余分块取法任意。
% \end{proof}

% \begin{example}
%     设 \(\varphi\) 是有限维线性空间 \(V\) 到 \(U\) 的线性映射，求证：必存在 \(U\) 到 \(V\) 的线性映射 \(\psi\)，使得 \(\varphi\psi\varphi = \varphi\)。
% \end{example}
    
% \begin{proof}[证法 2（代数方法）] \par
%     取定 \(V\) 和 \(U\) 的两组基，设 \(\varphi\) 在这两组基下的表示矩阵为 \(m \times n\) 矩阵 \(A\)，则由例\ref{ex3.93}可知，存在 \(n \times m\) 矩阵 \(B\)，使得 \(ABA = A\)。由矩阵 \(B\) 可定义从 \(U\) 到 \(V\) 的线性映射 \(\psi\)，它适合 \(\varphi\psi\varphi = \varphi\)。
% \end{proof}

% \begin{example}
%     设 \( V_1, V_2 \) 是 \( V \) 上线性变换 \( \varphi \) 的不变子空间，求证：\( V_1 \cap V_2, V_1 + V_2 \) 也是 \( \varphi \) 的不变子空间。
% \end{example}
    
% \begin{proof}
%     任取 \( v \in V_1 \cap V_2 \)，则由 \( v \in V_i \) 可得 \( \varphi(v) \in V_i \; (i = 1, 2) \)，于是 \( \varphi(v) \in V_1 \cap V_2 \)，从而 \( V_1 \cap V_2 \) 是 \( \varphi \) 的不变子空间。
    
%     任取 \( v \in V_1 + V_2 \)，则 \( v = v_1 + v_2 \)，其中 \( v_i \in V_i \)，故 \( \varphi(v_i) \in V_i \; (i = 1, 2) \)，于是
%     \[
%     \varphi(v) = \varphi(v_1) + \varphi(v_2) \in V_1 + V_2,
%     \]
%     从而 \( V_1 + V_2 \) 是 \( \varphi \) 的不变子空间。
% \end{proof}

\begin{example}
    用几何方法证明 Cayley-Hamilton定理：已知$\mathcal{A}  $是有限维线性空间V中的线性变换，\(f(\lambda)\) 是变换 \(\mathcal{A}\) 的特征多项式，则 \(f(\mathcal{A}) = 0\)。
\end{example}

\begin{proof}
    要证明 \(f(\mathcal{A}) = 0\)，只要证明 \(f(\mathcal{A}) \xi = 0\) 对任意 \(\xi \in V\) 都成立即可。 \par
    任取 $\xi  \in V$ 若$\xi = 0$ ，则$f(\mathcal{A}) \xi = 0$ 显然成立。 \par
    若$\xi \neq 0$ ，考虑$\mathcal{A}\xi$,若$\xi$与$\mathcal{A}\xi$ 线性相关，则停止，令$W = L(\xi)$，否则继续考虑$\mathcal{A}^2\xi$，若$\{\xi,\mathcal{A}\xi,\mathcal{A}^2\xi\}$ 线性相关，则停止，令$W = L(\xi,\mathcal{A}\xi)$，否则继续讨论下去。由于V是有限维线性空间，$dim V = n < \infty$，V中$k >n$个向量必定线性相关，上述讨论过程一定可以停止，所以不妨设存在正整数$m$使得$\{\xi,\mathcal{A}\xi,\cdots,\mathcal{A}^{m-1}\xi\}$ 线性无关，而$\{\xi,\mathcal{A}\xi,\cdots,\mathcal{A}^m\xi\}$ 线性相关。 \par
    所以$\mathcal{A}^m\xi$可由线性无关的向量组$\{\xi,\mathcal{A}\xi,\cdots,\mathcal{A}^{m-1}\xi\}$ 线性表示，即存在数 $a_0, a_1, \cdots, a_{m-1}$ 使得
\[
\mathcal{A}^m\xi = a_0 \xi + a_1 \mathcal{A} \xi + \cdots + a_{m-1} \mathcal{A}^{m-1} \xi.
\]
    令$W = L(\xi,\mathcal{A}\xi,\cdots,\mathcal{A}^{m-1}\xi)$，则W为不变子空间。将W中的一组基$\{\xi,\mathcal{A}\xi,\cdots,\mathcal{A}^{m-1}\xi\}$，扩张成V中的一组基$\{\xi,\mathcal{A}\xi,\cdots,\mathcal{A}^{m-1}\xi,v_{m+1},\cdots,v_{n}\}$则$\mathcal{A}$在这组基下的表示矩阵为
\[
\begin{pmatrix}
A_{11} & A_{12} \\
O & A_{22}
\end{pmatrix}
\]
其中$A_{11}$为线性变换$\mathcal{A}$限制在W上的表示矩阵。且
\[
A_{11} = \begin{pmatrix}
0 & 0 & \cdots & 0 & a_0 \\
1 & 0 & \cdots & 0 & a_1 \\
0 & 1 & \cdots & 0 & a_2 \\
\vdots & \vdots & \ddots & \vdots & \vdots \\
0 & 0 & \cdots & 1 & a_{m-1}
\end{pmatrix}
\]

\[
f(\lambda) = \begin{vmatrix}
\lambda I_n - \begin{pmatrix}
A_{11} & A_{12} \\
O & A_{22}
\end{pmatrix}
\end{vmatrix} = \begin{vmatrix}
\lambda I_m - A_{11} & -A_{12} \\
O & \lambda I_{n-m} - A_{22}
\end{vmatrix} = \begin{vmatrix}
\lambda I_m - A_{11}
\end{vmatrix} \begin{vmatrix}
\lambda I_{n-m} - A_{22}
\end{vmatrix}
\]
令$g(\lambda) = \begin{vmatrix}
    \lambda I_m - A_{11}
    \end{vmatrix}$，
$h(\lambda) =  \begin{vmatrix}
    \lambda I_{n-m} - A_{22}
    \end{vmatrix}$

\[
g(\lambda) = \begin{vmatrix}
\lambda & 0 & \cdots & 0 & -a_0 \\
-1 & \lambda & \cdots & 0 & -a_1 \\
0 & -1 & \cdots & 0 & -a_2 \\
\vdots & \vdots & \ddots & \vdots & \vdots \\
0 & 0 & \cdots & -1 & \lambda-a_{m-1}
\end{vmatrix} = \lambda^m - a_{m-1} \lambda^{m-1} - \cdots - a_1 \lambda - a_0
\]

   从而$g(\mathcal{A}) \xi = \mathcal{A}^m\xi - a_{m-1} \mathcal{A}^{m-1}\xi - \cdots - a_1 \mathcal{A}\xi - a_0 \xi = 0 $ (利用的是$\mathcal{A}^m\xi$的线性表示)

   由于同一变换的多项式具有可交换性，故
\[
f(\mathcal{A}) \xi = g(\mathcal{A}) h(\mathcal{A}) \xi = h(\mathcal{A}) g(\mathcal{A}) \xi  = h(\mathcal{A})0 = 0.
\]
    因此 $f(\mathcal{A}) = 0$。

\end{proof}



\chapter{第二次习题课}

\section{不定积分}

\begin{example}
    求下列不定积分：
    \[
    \int \frac{\ln x}{x\sqrt{1 + \ln x}} \; dx
    \]
    \[
    \int (1 + x - \frac{1}{x})e^{x + \frac{1}{x}} \; dx
    \]
    \[
    \int \frac{1}{(x+1)(x^2 - 2x + 2)} \; dx
    \]
    \[
    \int \frac{x}{1 + \sin x} \; dx
    \]
    \[
    \int \frac{1}{1 + x^3} \; dx
    \]
    \[
    \int \frac{\sin x}{1 + \sin 2x} \; dx
    \]
    \[
    \int \frac{1}{1 + x^5} \; dx
    \]
\end{example}

\begin{example}
    \begin{enumerate}
        \item 用配对积分法计算下列不定积分：
        \begin{enumerate}
            \item \(\displaystyle \int \frac{e^x \, dx}{e^x + e^{-x}}\);
            \item \(\displaystyle \int \frac{b \sin x + a \cos x}{a \sin x + b \cos x} \, dx \quad (a \neq b)\);
        \end{enumerate}
        
        % \item 通过计算不定积分
        % \[
        % \int (\cos^4 x - \sin^4 x) \, dx \quad \text{与} \quad \int (\cos^4 x + \sin^4 x) \, dx,
        % \]
        \item 
        求出不定积分
        \[
        \int \cos^4 x \, dx \quad \text{和} \quad \int \sin^4 x \, dx.
        \]
    \end{enumerate}
\end{example}

\begin{solution}
(a)
令 \( M(x) = \int \frac{e^x}{e^x + e^{-x}} \, dx \), \( N(x) = \int \frac{e^{-x}}{e^x + e^{-x}} \, dx \)，则
\[
M(x) + N(x) = x + c,
\]
\[
M(x) - N(x) = \ln(e^x + e^{-x}) + c,
\]
故
\[
M(x) = \frac{1}{2}x + \frac{1}{2} \ln(e^x + e^{-x}) + c.
\]

(b)
令 \( M(x) = \int \frac{\sin x}{a \sin x + b \cos x} \, dx \), \( N(x) = \int \frac{\cos x}{a \sin x + b \cos x} \, dx \)，则
\[
aM(x) + bN(x) = x + c,
\]
\[
aN(x) - bM(x) = \ln|a \sin x + b \cos x| + c,
\]
故原积分
\[
\begin{aligned}
&\int \frac{b \sin x + a \cos x}{a \sin x + b \cos x} \, dx \\
&= bM(x) + aN(x) \\
&= \frac{b}{a^2+b^2}(ax - b\ln|a\sin x + b\cos x|) + \frac{a}{a^2+b^2}(bx + a\ln|a\sin x + b\cos x|) + c \\
&= \frac{2ab}{a^2+b^2}\,x - \frac{a^2-b^2}{a^2+b^2}\ln|a\sin x + b\cos x| + c.
\end{aligned}
\]
\end{solution}

\begin{solution}
解1
用分部积分法可进行如下：
\[
\begin{aligned}
I &= -\int \sin^3 x \, dx \\
    &= -\sin^3 x \cos x + \int \cos x \, d\sin^3 x \\
    &= -\sin^3 x \cos x + 3 \int \sin^2 x \cos^2 x \, dx \\
    &= -\sin^3 x \cos x + 3 \int \sin^4 x (1 - \sin^2 x) \, dx \\
    &= -\sin^3 x \cos x + 3 \int \sin^2 x \, dx - 3 \int \sin^4 x \, dx \\
    &= -\sin^3 x \cos x + \frac{3}{2} \int (1 - \cos 2x) \, dx - 3I \\
    &= -\sin^3 x \cos x + \frac{3}{2} x - \frac{3}{8} \sin 2x - 3I,
\end{aligned}
\]
所以得到
\[
I = -\frac{1}{4} \sin^3 x \cos x + \frac{3}{8} x - \frac{3}{16} \sin 2x + C.
\]

解2
如果用三角函数的倍角公式，则可得到更简单的解法：
\[
\begin{aligned}
I &= \int \left( \frac{1 - \cos 2x}{2} \right)^2 \, dx = \frac{1}{4} \int (1 - 2\cos 2x + \cos^2 2x) \, dx \\
    &= \frac{1}{4} \int \left( 1 - 2\cos 2x + \frac{1 + \cos 4x}{2} \right) \, dx \\
    &= \frac{3}{8} x - \frac{1}{4} \sin 2x + \frac{1}{32} \sin 4x + C.
\end{aligned}
\]

解3
\[
\begin{aligned}
\int (\cos^4 x - \sin^4 x) \, dx &= \int (\cos^2 x - \sin^2 x) \, dx = \sin 2x + C, \\
\int (\cos^4 x + \sin^4 x) \, dx &= \int (1 - 2\sin^2 x \cos^2 x) \, dx = \int \left(1 - \frac{1}{2}\sin^2 2x\right) dx \\
&= \frac{3}{4}x - \frac{1}{4}\sin 4x + C.
\end{aligned}
\]
\end{solution}

\begin{example}
    计算下列不定积分：
    \begin{enumerate}
        \item \[
    \int \frac{dx}{1 + x^4}.
    \]
        \item \[
    \int \frac{1}{1 + \tan x} \; dx
    \]
    \end{enumerate}
    
\end{example}

\begin{solution}
    令 \( M(x) = \int \frac{dx}{1+x^4}, N(x) = \int \frac{x^2 dx}{1+x^4} \)，则有
    \[
    \begin{aligned}
    M(x) - N(x) &= \int \frac{1-x^2}{1+x^4} dx = -\int \frac{1-\frac{1}{x^2}}{x^2+\frac{1}{x^2}} dx \\
    &= -\int \frac{d\left(x+\frac{1}{x}\right)}{\left(x+\frac{1}{x}\right)^2-2} = -\frac{1}{2\sqrt{2}}\ln\frac{x^2-\sqrt{2}x+1}{x^2+\sqrt{2}x+1} + C, \\
    M(x) + N(x) &= \int \frac{1+x^2}{1+x^4} dx = \int \frac{1+\frac{1}{x^2}}{x^2+\frac{1}{x^2}} dx = \int \frac{d\left(x-\frac{1}{x}\right)}{\left(x-\frac{1}{x}\right)^2+2} \\
    &= \frac{1}{\sqrt{2}}\arctan\frac{x-\frac{1}{x}}{\sqrt{2}} + C = \frac{1}{\sqrt{2}}\arctan\frac{x^2-1}{\sqrt{2}x} + C,
    \end{aligned}
    \]
    因此得到
    \[
    \begin{aligned}
    M(x) &= \frac{1}{2}\bigl[(M(x) + N(x)) + (M(x) - N(x))\bigr] \\
    &= -\frac{1}{4\sqrt{2}}\ln\frac{x^2-\sqrt{2}x+1}{x^2+\sqrt{2}x+1} + \frac{1}{2\sqrt{2}}\arctan\frac{x^2-1}{\sqrt{2}x} + C.
    \end{aligned}
    \]
    本题也可用因式分解裂项的标准方法做，但计算量却要大得多。
\end{solution}

\begin{solution}
    参考上例组合积分 
    \[
    \begin{aligned}
        \int \frac{1}{1 + \tan x} \; dx = \int \frac{\cos x}{\cos x + \sin x} \; dx \\
        = \int \frac{\frac{\cos x + \sin x  + \cos x- \sin x }{2}}{\cos x + \sin x} \; dx\\
        = \int \frac{1}{2} \; dx + \int \frac{\cos x - \sin x}{2(\cos x + \sin x)} \; dx\\
        = \frac{1}{2}x + \frac{1}{2}\ln|\cos x + \sin x| + C.
    \end{aligned}
    \]
\end{solution}

\begin{example}
    设对正整数 \( n > 2 \)，定义 
    \[
    I_n = \int \frac{\sin nx}{\sin x} \, dx.
    \]
    证明：
    \[
    I_n = \frac{2}{n-1} \sin(n-1)x + I_{n-2}.
    \]
    \end{example}
    
\begin{proof}
考虑降 \( n \)，
\[
\begin{aligned}
I_n &= \int \frac{\sin(n-1)x \cos x + \sin x \cos(n-1)x}{\sin x} \, dx \\
&= \int \frac{\sin(n-1)x \cos x}{\sin x} \, dx + \int \cos(n-1)x \, dx \\
&= \frac{1}{2} \int \frac{\sin nx + \sin(n-2)x}{\sin x} \, dx + \int \cos(n-1)x \, dx \\
&= \frac{1}{2} I_n + \frac{1}{2} I_{n-2} + \frac{1}{n-1} \sin(n-1)x,
\end{aligned}
\]
所以
\[
I_n = \frac{2}{n-1} \sin(n-1)x + I_{n-2}.
\]
\end{proof}

\section{习题讲解}

\begin{example}
    若$\lambda_1,\lambda_2,\cdots,\lambda_k$是线性变换$\mathcal{A}$的k个互不相同的特征值，而$\{\xi_{i1},\xi_{i2},\cdots,\xi_{ir_i}\} (i = 1 ,2,\cdots,k)$是属于特征值$\lambda_i$的线性无关的特征向量组。求证：向量组$\{\xi_{11},\xi_{12},\cdots,\xi_{1r_1},\cdots,\xi_{k1},\xi_{k2},\cdots,\xi_{kr_k}\}$也线性无关。
\end{example}

\begin{proof}
    数学归纳法：当$k = 1$时，结论显然成立。假设当$k = m-1$时结论成立，现讨论$k = m$的情况。设有数$a_{ij} (1 \leq i \leq m, 1 \leq j \leq r_i)$使得
\[a_{11}\xi_{11} + a_{12}\xi_{12} + \cdots + a_{1r_1}\xi_{1r_1} + \cdots + a_{m1}\xi_{m1} + a_{m2}\xi_{m2} + \cdots + a_{mr_m}\xi_{mr_m} = 0\]

    上式两边同时作用$\mathcal{A}$，则有
\[
a_{11}\lambda_1\xi_{11} + a_{12}\lambda_1\xi_{12} + \cdots + a_{1r_1}\lambda_1\xi_{1r_1} + \cdots + a_{m1}\lambda_m\xi_{m1} + a_{m2}\lambda_m\xi_{m2} + \cdots + a_{mr_m}\lambda_m\xi_{mr_m} = 0
\]

    上式两边同时乘以$\lambda_m$，则有
\[
a_{11}\lambda_m\xi_{11} + a_{12}\lambda_m\xi_{12} + \cdots + a_{1r_1}\lambda_m\xi_{1r_1} + \cdots + a_{m1}\lambda_m\xi_{m1} + a_{m2}\lambda_m\xi_{m2} + \cdots + a_{mr_m}\lambda_m\xi_{mr_m} = 0
\]
    上式两边相减，得到

\(
a_{11}(\lambda_1 - \lambda_m)\xi_{11} + a_{12}(\lambda_1 - \lambda_m)\xi_{12} + \cdots + a_{1r_1}(\lambda_1 - \lambda_m)\xi_{1r_1} + \cdots + a_{m-1,1}(\lambda_{m-1} - \lambda_m)\xi_{m-1,1} +  a_{m-1,2}(\lambda_{m-1} - \lambda_m)\xi_{m-1,2} + \cdots + a_{m-1,r_{m-1}}(\lambda_{m-1} - \lambda_m)\xi_{m-1,r_{m-1}} = 0
\)

    由归纳假设可知，$a_{ij}(\lambda_i - \lambda_m) = 0$，由于$\lambda_i \neq \lambda_m$，故$a_{ij} = 0$，从而$a_{m1}\xi_{m1} + a_{m2}\xi_{m2} + \cdots + a_{mr_m}\xi_{mr_m} = 0$，又由于$\{\xi_{m1},\xi_{m2},\cdots,\xi_{mr_m}\}$线性无关，所以$a_{m1} = a_{m2} = \cdots = a_{mr_m} = 0$。
    
    因此$a_{ij} = 0(1 \leq i \leq m, 1 \leq j \leq r_i)$，结论成立。
\end{proof}

\begin{example}\label{ex6.1.2}
    设 \( A \) 是 \( n \) 阶复方阵，证明：存在 \( n \) 阶非异阵 \( P \)，使得 \( P^{-1}AP \) 是一个上三角阵。
\end{example}

\begin{proof}
    设 \( A \) 是 \( n \) 阶复方阵，现对 \( n \) 用数学归纳法。当 \( n = 1 \) 时结论显然成立。假设对 \( n - 1 \) 阶矩阵结论成立，现对 \( n \) 阶矩阵 \( A \) 来证明。设 \( \lambda_1 \) 是 \( A \) 的一个特征值，则存在非零列向量 \( \alpha_1 \)，使
    \[
    A\alpha_1 = \lambda_1 \alpha_1.
    \]
    将 \( \alpha_1 \) 作为 \( \mathbb{C}^n \) 的一个基向量，并扩展为 \( \mathbb{C}^n \) 的一组基 \( \{ \alpha_1, \alpha_2, \cdots, \alpha_n \} \)。将这些基向量按照列分块方式拼成矩阵 \( P = (\alpha_1, \alpha_2, \cdots, \alpha_n) \)，则 \( P \) 为 \( n \) 阶非异阵，且
    \[
    AP = A(\alpha_1, \alpha_2, \cdots, \alpha_n) = (A\alpha_1, A\alpha_2, \cdots, A\alpha_n)
    = (\alpha_1, \alpha_2, \cdots, \alpha_n)
    \begin{pmatrix}
    \lambda_1 & * \\
    0 & A_1
    \end{pmatrix},
    \]
    其中 \( A_1 \) 是一个 \( n - 1 \) 阶方阵。注意到 \( P = (\alpha_1, \alpha_2, \cdots, \alpha_n) \) 非异，上式即为
    \[
    P^{-1}AP = 
    \begin{pmatrix}
    \lambda_1 & * \\
    0 & A_1
    \end{pmatrix}.
    \]
    因为 \( A_1 \) 是一个 \( n - 1 \) 阶方阵，所以由归纳假设可知，存在 \( n - 1 \) 阶非异阵 \( Q \)，使 \( Q^{-1}A_1Q \) 是一个上三角阵。令
    \[
    R = 
    \begin{pmatrix}
    1 & O \\
    O & Q
    \end{pmatrix},
    \]
    则 \( R \) 是 \( n \) 阶非异阵，且
    \[
    \begin{aligned}
    R^{-1}P^{-1}APR &=
    \begin{pmatrix}
    1 & O \\
    O & Q
    \end{pmatrix}^{-1}
    \begin{pmatrix}
    \lambda_1 & * \\
    0 & A_1
    \end{pmatrix}
    \begin{pmatrix}
    1 & O \\
    O & Q
    \end{pmatrix} \\
    &=
    \begin{pmatrix}
    1 & O \\
    O & Q^{-1}
    \end{pmatrix}
    \begin{pmatrix}
    \lambda_1 & * \\
    0 & A_1
    \end{pmatrix}
    \begin{pmatrix}
    1 & O \\
    O & Q
    \end{pmatrix} \\
    &=
    \begin{pmatrix}
    \lambda_1 & * \\
    O & Q^{-1}A_1Q
    \end{pmatrix}.
    \end{aligned}
    \]
    这是一个上三角阵，它与 \( A \) 相似。
\end{proof}

\begin{example}
    用相似变换代数方法证明Cayley-Hamilton 定理：设 \( A \) 是数域 \( K \) 上的 \( n \) 阶矩阵，\( f(x) \) 是 \( A \) 的特征多项式，证明：\( f(A) = O \)。
\end{example}

\begin{proof}
    由例 \ref{ex6.1.2} 知 \( A \) 复相似于一个上三角阵，也就是说存在可逆阵 \( P \)，使
    \[
    P^{-1}AP = B
    \]
    是一个上三角阵，其中 \( P \) 与 \( B \) 都是复矩阵，但 \( A \) 与 \( B \) 有相同的特征多项式 \( f(x) \)。记
    \[
    f(x) = x^n + a_1x^{n-1} + \cdots + a_n,
    \]
    则 \( f(B) = O \)。而
    \[
\begin{aligned}
    f(A) &= A^n + a_1A^{n-1} + \cdots + a_nI_n \\
    &= (PBP^{-1})^n + a_1(PBP^{-1})^{n-1} + \cdots + a_nI_n \\
    &= PB^nP^{-1} + a_1PB^{n-1}P^{-1} + \cdots + a_nI_n \\
    &= P(B^n + a_1B^{n-1} + \cdots + a_nI_n)P^{-1} \\
    &= Pf(B)P^{-1} = O.
\end{aligned}
    \]
\end{proof}

\begin{example}\label{ex6.7}
    设 \( n \) 阶矩阵 \( A \) 的全体特征值为 \( \lambda_1, \lambda_2, \cdots, \lambda_n \)，\( f(x) \) 是一个多项式，求证：\( f(A) \) 的全体特征值为 \( f(\lambda_1), f(\lambda_2), \cdots, f(\lambda_n) \)。
\end{example}
    
\begin{proof}
    因为任一 \( n \) 阶矩阵均复相似于上三角矩阵，故可设
    \[
    P^{-1}AP = 
    \begin{pmatrix}
    \lambda_1 & * & * \\
    0 & \lambda_2 & * \\
    \vdots & \vdots & \ddots \\
    0 & 0 & \cdots & \lambda_n
    \end{pmatrix}.
    \]
    注意到上三角矩阵的和、数乘及乘方仍是上三角矩阵，经计算可得
    \[
    P^{-1}f(A)P = f(P^{-1}AP) =
    \begin{pmatrix}
    f(\lambda_1) & * & * \\
    0 & f(\lambda_2) & * \\
    \vdots & \vdots & \ddots \\
    0 & 0 & \cdots & f(\lambda_n)
    \end{pmatrix},
    \]
    因此 \( f(A) \) 的全体特征值为 \( f(\lambda_1), f(\lambda_2), \cdots, f(\lambda_n) \)。
\end{proof}
    
\begin{remark}
    例 \ref{ex6.7} 告诉我们：如果能将一个复杂矩阵写成一个简单矩阵的多项式，那么就可由简单矩阵的特征值得到复杂矩阵的特征值。
\end{remark}

\begin{example}
    设 \( n \) 阶矩阵 \( A \) 适合一个多项式 \( g(x) \)，即 \( g(A) = O \)。求证：\( A \) 的特征值 \( \lambda_0 \) 也适合 \( g(x) \)，即 \( g(\lambda_0) = 0 \)。
\end{example}
    
\begin{proof}
    由例 \ref{ex6.7}可知，\( g(\lambda_0) \) 是 \( g(A) = O \) 的特征值，从而 \( g(\lambda_0) = 0 \)。
\end{proof}

\begin{example}
    
\end{example}

\begin{proof}
    
\end{proof}



% \nocite{*}
% \printbibliography[heading=bibintoc, title=\ebibname]
% \appendix
% \chapter{基本数学工具}
% 本附录包括了计量经济学中用到的一些基本数学，我们扼要论述了求和算子的各种性质，研究了线性和某些非线性方程的性质，并复习了比例和百分数。我们还介绍了一些在应用计量经济学中常见的特殊函数，包括二次函数和自然对数，前 4 节只要求基本的代数技巧，第 5 节则对微分学进行了简要回顾；虽然要理解本书的大部分内容，微积分并非必需，但在一些章末附录和第 3 篇某些高深专题中，我们还是用到了微积分。

% \section{求和算子与描述统计量}

% \textbf{求和算子} 是用以表达多个数求和运算的一个缩略符号，它在统计学和计量经济学分析中扮演着重要作用。如果 $\{x_i: i=1, 2, \ldots, n\}$ 表示 $n$ 个数的一个序列，那么我们就把这 $n$ 个数的和写为：

% \begin{equation}
% \sum_{i=1}^n x_i \equiv x_1 + x_2 +\cdots + x_n
% \end{equation}

\end{document}